\documentclass[12pt, a4paper]{article}

\usepackage[utf8]{inputenc}
\usepackage[T1]{fontenc}
\usepackage{geometry}
\usepackage{xcolor}
\usepackage{titlesec}
\usepackage{fancyhdr}
\usepackage{enumitem}
\usepackage{hyperref}
\usepackage{parskip}
\usepackage{amsmath, amssymb}
\usepackage{tcolorbox}
\tcbuselibrary{skins,breakable}

\geometry{top=2.5cm,bottom=2.5cm,left=2.5cm,right=2.5cm,headheight=15pt}

\definecolor{azul42}{HTML}{1F3A8A}
\definecolor{azulclaro}{HTML}{E8EEF9}
\definecolor{cinza}{HTML}{555555}

\hypersetup{colorlinks=true,linkcolor=azul42,urlcolor=azul42}

\pagestyle{fancy}
\fancyhf{}
\fancyhead[L]{\color{azul42}\small\textit{Clube de Matemática @ 42 São Paulo}}
\fancyhead[R]{\color{cinza}\small\textit{Reunião \#000}}
\fancyfoot[C]{\color{cinza}\small\thepage}
\renewcommand{\headrulewidth}{0.4pt}
\renewcommand{\headrule}{\color{azul42}\hrule width\headwidth height\headrulewidth}
\renewcommand{\footrulewidth}{0pt}

\titleformat{\section}
  {\color{azul42}\Large\bfseries}{\thesection.}{0.6em}{}
  [\vspace{2pt}{\color{azul42}\titlerule[1.5pt]}\vspace{4pt}]
\titleformat{\subsection}{\color{azul42}\large\bfseries}{\thesubsection}{0.6em}{}
\titlespacing{\section}{0pt}{18pt}{8pt}
\titlespacing{\subsection}{0pt}{12pt}{6pt}

\setlist[itemize]{leftmargin=1.5em,itemsep=3pt,parsep=0pt,topsep=4pt,
  label={\color{azul42}\textbullet}}

\newtcolorbox{infobox}[1][]{enhanced,breakable,
  colback=azulclaro,colframe=azul42,boxrule=0.8pt,arc=4pt,
  left=8pt,right=8pt,top=6pt,bottom=6pt,#1}

\newtcolorbox{citacao}{enhanced,breakable,
  colback=white,colframe=azul42!30,boxrule=0pt,leftrule=3pt,arc=0pt,
  left=10pt,right=8pt,top=4pt,bottom=4pt,
  fontupper=\itshape\color{cinza}}

\begin{document}
\thispagestyle{empty}

\begin{center}
  \vspace*{2cm}
  {\fontsize{48}{48}\selectfont\color{azul42}$\Sigma$}\\[0.6cm]
  {\fontsize{26}{30}\selectfont\color{azul42}\textbf{Reunião \#000}}\\[0.3cm]
  {\fontsize{14}{18}\selectfont\color{cinza}Sessão Teste --- Clube de Matemática @ 42 São Paulo}\\[0.8cm]
  {\color{azul42}\rule{0.5\textwidth}{1pt}}\\[0.6cm]
  {\normalsize\color{cinza}30 de abril de 2026}
\end{center}

\vspace{1.5cm}

\begin{infobox}
  \textbf{Formato:} Exploração coletiva sem cronômetro\\
  \textbf{Tema:} Questão proposta pelo grupo --- generalização de conversões de base numérica\\
  \textbf{Referências externas consultadas:} Nenhuma
\end{infobox}

\newpage

\section{Apresentações}

A sessão começou com cada membro apresentando seu nível de contato com a matemática --- do ensino médio a graduações em exatas. Não houve hierarquia nisso: o objetivo era entender de onde cada um parte, não classificar quem sabe mais.

Essa diversidade foi imediatamente útil. Ao longo da sessão, quem tinha mais base ajudava naturalmente quem tinha menos, sem que isso precisasse ser organizado.

\section{O Problema}

O problema da sessão foi uma representação matemática de \textit{Bit Masking} --- tema diretamente ligado ao cursus da 42.

Seja $U$ um conjunto universal finito com cinco variáveis:
$$U = \{x_0,\, x_1,\, x_2,\, x_3,\, x_4\}$$

Seja $S$ qualquer subconjunto de $U$ (ou seja, $S \subseteq U$). Definimos uma função característica $f(S)$ que mapeia um conjunto $S$ para uma string binária de 5 bits.

A função avalia os elementos de $S$ usando uma condição lógica para gerar cada bit $b_i$ da string $b_4b_3b_2b_1b_0$:
$$
b_i =
\begin{cases}
    1 & \text{se } x_i \in S \\
    0 & \text{se } x_i \notin S
\end{cases}
$$

\noindent(Nota: $b_0$ representa o bit menos significativo, à direita, correspondente a $x_0$.)

Dados os subconjuntos:
$$A = \{x_0,\, x_2,\, x_3\} \qquad B = \{x_1,\, x_2,\, x_4\}$$

\begin{enumerate}
  \item Qual é $f(A)$?
  \item Qual é $f(B)$?
\end{enumerate}

\section{Como a Sessão Aconteceu}

Ninguém sabia a resposta de antemão. O que aconteceu foi uma \textbf{pesquisa matemática coletiva e lenta} --- no bom sentido. Cada pessoa precisava externalizar seus pensamentos, suas ideias e, principalmente, suas dúvidas.

\begin{citacao}
``Nah, eu não sabia qual símbolo usar'' --- então só escreve por extenso mesmo, tipo: ``3 está no conjunto dos naturais''. A ideia é externalizar o que sabe, para assim os outros poderem aprender, ajudar e acompanhar.
\end{citacao}

Quando a notação faltava, o grupo simplesmente inventava a própria. A formalização veio depois, como abstração do que já havia sido pensado --- da mesma forma que a matemática foi construída historicamente.

Não houve professor. Não houve cronômetro. O modelo foi o da \textbf{Grécia antiga}: pensar em voz alta, divergir, questionar e reconstruir o raciocínio coletivamente.

\section{O Raciocínio}

O raciocínio foi construído em três etapas, cada uma surgindo naturalmente da anterior.

\subsection{Etapa 1 --- Representação Binária (Herom)}

Herom começou desenhando e indagando: como representar formalmente se um elemento pertence ou não a um conjunto? A resposta imediata foi a função característica binária, que o grupo escreveu assim:

$$
b_i = \begin{cases} 1 & \text{se } x_i \in S \\ 0 & \text{se } x_i \notin S \end{cases}
$$

Esse é exatamente o $f(S)$ do problema. Cada bit $b_i$ registra a presença ou ausência de $x_i$ no subconjunto $S$.

\subsection{Etapa 2 --- Generalização para Base 3 (Breno)}

A primeira ideia de generalização veio do Breno: e se o conjunto $S$ fosse \textit{ordenado}? Um elemento poderia estar presente de duas formas distintas --- ordenado ou desordenado. Isso naturalmente pede uma terceira condição, e portanto uma base 3:

$$
T_i = \begin{cases} 0 & \text{se } x_i \notin S \\ 1 & \text{se } x_i \in S \text{ e não está ordenado} \\ 2 & \text{se } x_i \in S \text{ e está ordenado } \end{cases}
$$

Com isso, a string resultante deixa de ser binária e passa a ser ternária --- cada dígito carrega mais informação sobre a estrutura do conjunto.

\subsection{Etapa 3 --- Generalização para Base $m$ (Felipe / \texttt{atoi\_base})}

A partir da ideia do \texttt{atoi\_base} trazida por Felipe --- função que converte strings para inteiros em bases arbitrárias --- o grupo percebeu que as etapas anteriores eram casos particulares de uma função mais geral.

Para qualquer base $m$, define-se $R_i$ como:

$$
R_i: \begin{cases}
  n = \dfrac{x_i}{m} \\[6pt]
  n_i = 0 & \text{se } x_i \notin S \\[4pt]
  n_i = 1 & \text{se } x_i \in S \text{ (sendo } m \text{ a base de } S\text{)} \\[4pt]
  n_j = \dfrac{x_j}{m} & \text{se } x_i \text{ for o último elemento}
\end{cases}
$$

O exemplo produzido no quadro foi $\{x_1, x_0, x_3\}$, aplicando $R_i$ sobre esse subconjunto ordenado.

A conexão com \texttt{atoi\_base} foi o que permitiu enxergar a estrutura geral: assim como \texttt{atoi} reconstrói um número a partir de seus dígitos em uma base qualquer, $R_i$ constrói uma representação posicional a partir da pertinência e ordem dos elementos em $S$.

\begin{infobox}
  O problema, a generalização para base 3 e a função $R_i$ foram inteiramente inventados pelo grupo durante a sessão. Nenhuma referência externa foi consultada.
\end{infobox}

\section{Observações sobre o Formato}

Esta primeira sessão também serviu para apresentar o Estatuto do clube e discutir sua filosofia. Alguns pontos que emergiram na prática:

\begin{itemize}
  \item Pessoas sem base matemática formal conseguem fazer perguntas que avançam a discussão tanto quanto quem tem;
  \item A ausência de pesquisa durante a sessão não foi uma limitação --- foi o que tornou o pensamento genuinamente coletivo;
  \item Inventar notação antes de formalizá-la revelou \textit{por que} a notação matemática existe do jeito que existe.
\end{itemize}

\begin{citacao}
A ideia é aprender do ``por que'' a matemática foi desenhada do jeito que é hoje --- descobrindo, igual como os matemáticos antigamente faziam.
\end{citacao}

\section{Próximos Passos}

A ideia que ficou: cada membro deve mergulhar em textos, livros e referências \textbf{antes} da próxima sessão --- não para chegar com a resposta, mas para chegar com pensamentos, dúvidas e ideias que alimentem a discussão.

O tema da próxima sessão emergirá do backlog coletivo.

\end{document}
